\section{Results} \label{sec:results}

\subsection{Financial Proficiency, Savings and Consumption Attitudes, and Behavioral Outcomes} \label{subsec:literacy}

%In this section, we present ITT estimates of the Financial Literacy pilot intervention on students' financial proficiency, attitudes towards hypothetical settings involving savings and consumption decisions, and behavioral outcomes. We show the effects for the pooled sample, for elementary and middle school grades, and for each grade separately.  

\autoref{tab:ITT} shows that the intervention has a positive and statistically significant impact of 0.07 of a standard deviation on students' financial proficiency, with the effect driven by the middle school grades. Larger impacts on students of more advanced grades are consistent with the hypothesis that financial education programs are more effective in increasing the financial proficiency of the students more likely to experience financial decisions in daily life (\cite{frisancho2020impact}, \cite{Bruhn2016}), and \cite{zia2023financial}. Our results are in line with \cite{jamison2014financial} as the authors find an impact of 0.08 of a standard deviation of a financial literacy pilot implemented among Ugandan youth groups, but are considerably smaller when compared to the ones reported by \cite{de2021effect} and \cite{batty2020experiential} who find an increase in the financial literacy of middle and elementary grade students of 0.46 and 0.25 of a standard deviation, respectively. 

Our estimates also point to positive impacts on consumption and saving indices, the ones measuring students' attitudes, particularly among fifth-graders (\autoref{tab:ITT}).\footnote{Due to characteristics of the research instruments for the third graders, the variables employed to set up the savings and consumption index are unavailable for this grade. Thus, column (2) estimates correspond to only fifth-graders.} These estimates should be interpreted as children becoming more conscious, self-controlled, and forward-looking. Interestingly, the magnitude of the impact on saving attitudes declines steadily with school grades. This pattern suggests that older students are less risk-averse and potentially self-controlled. 

The larger point estimates for self-control and risk aversion among fifth graders might be associated with a few factors. Younger students might be less overconfident than teenagers and this can make them more responsive to interventions aimed to develop risk-aversion behavior. At the same token, it might be easier to change the behavior of younger students since their financial habits are not consolidated yet.  

%\newgeometry{left = 1.5cm, right = 4cm, top = 2cm, bottom = -2cm}

\vspace{1cm}
\begin{table}[h!]
\centering
\fontsize{9}{10} \selectfont
\caption{ITT estimates for financial proficiency and hypothetical attitudes towards savings and consumption}
    \begin{tabular}{lccccccc}
    \toprule
    \toprule
          & (1)   & (2)   & (3)   & (4)   & (5)   & (6)   & (7) \\
          & Pooled & Elementary  & Middle school & 3\textsuperscript{rd} grade & 5\textsuperscript{th} grade & 7\textsuperscript{th} grade & 9\textsuperscript{th} grade \\
    \midrule
    \textbf{Financial Proficiency} &       &       &       &       &       &       &  \\
   
 Treatment &    0.073 &      0.047 &    0.099 &      0.053 &      0.043 &     0.090 &     0.106 \\

           &    (0.034) &    (0.046) &    (0.042) &    (0.077) &    (0.046) &    (0.048) &    (0.055) \\

    pvalue &      0.033 &      0.307 &      0.022 &      0.495 &       0.35 &      0.061 &      0.057 \\

 pvalue RI &      0.034 &      0.328 &      0.024 &      0.502 &      0.381 &      0.055 &      0.064 \\

    N. obs &      14655 &       7641 &       7014 &       3635 &       4006 &       3640 &       3374 \\

Number of clusters &        201 &        172 &        145 &        170 &        168 &        144 &        144 \\

        R2 &      0.167 &       0.15 &      0.188 &      0.127 &      0.173 &      0.208 &      0.168 \\


    Strata F.E. & Yes   & Yes   & Yes   & Yes   & Yes   & Yes   & Yes \\
          &       &       &       &       &       &       &  \\
    \midrule
    \textbf{Consumption index} &       &       &       &       &       &       &  \\
        \textbf{self-control} &       &       &       &       &       &       &  \\
 Treatment &   0.094 &   0.122 &   0.082 &            &   0.122 &     0.080 &    0.084 \\

           &    (0.026) &    (0.042) &    (0.031) &            &    (0.042) &    (0.041) &    (0.037) \\

    pvalue &          0 &      0.004 &       0.01 &            &      0.004 &      0.054 &      0.025 \\

 pvalue RI &      0.002 &      0.006 &      0.009 &            &      0.009 &      0.052 &      0.029 \\

    N. obs &      10220 &       3648 &       6572 &            &       3648 &       3409 &       3163 \\

Number of clusters &        200 &        167 &        145 &            &        167 &        144 &        142 \\

        R2 &      0.089 &      0.119 &      0.076 &            &      0.119 &        0.1 &      0.055 \\

    Strata F.E. & Yes   & Yes   & Yes   & -     & Yes   & Yes   & Yes \\
          &       &       &       &       &       &       &  \\
    \midrule
    \textbf{Saving index} &       &       &       &       &       &       &  \\
      \textbf{risk-aversion} &       &       &       &       &       &       &  \\

 Treatment &    0.055 &    0.087 &      0.044 &            &    0.087 &      0.066 &      0.018 \\

           &    (0.025) &    (0.041) &    (0.029) &            &    (0.041) &    (0.040) &    (0.038) \\

    pvalue &      0.031 &      0.036 &      0.138 &            &      0.036 &      0.102 &      0.644 \\

 pvalue RI &      0.027 &      0.042 &      0.114 &            &      0.042 &      0.105 &      0.645 \\

    N. obs &      10022 &       3549 &       6473 &            &       3549 &       3368 &       3105 \\

Number of clusters &        200 &        166 &        145 &            &        166 &        144 &        142 \\

        R2 &      0.046 &      0.096 &      0.025 &            &      0.096 &      0.032 &      0.019 \\

    Strata F.E. & Yes   & Yes   & Yes   & -     & Yes   & Yes   & Yes \\
    \bottomrule
    \end{tabular}%    
\label{tab:ITT}

\begin{tablenotes}
\fontsize{9}{10} \selectfont
    \item Note: Clustered standard errors at the school level in parenthesis. RI stands for Randomization Inference. \\
      The financial proficiency is computed by CAEd based on students' answers on a standardized exam that uses the Item Response Theory (IRT). The consumption and saving indices are based on students’ choices in questions containing four possible answers (I totally agree, I agree, I disagree, I totally disagree) and they are aimed to measure students' self-control and risk-aversion, respectively. The financial proficiency, savings, and consumption indices are normalized so that the treatment effect can be interpreted in terms of standard deviations (SD). 
\end{tablenotes}
\end{table}%



To investigate whether the intervention had distributional impacts on financial proficiency, we estimate the quantile treatment effects.\footnote{Since the estimates are on the eligible population and do not control for covariates, we call it unconditional quantile intent to treat effects (UQITT)}. \autoref{fig:te_firstcycle} displays the effect of the pilot on elementary school students at 20 different quantiles of this distribution. The UQITT estimates show a positive and statistically significant of about 0.05 of an SD starting in the 45 percentile of the financial proficiency distribution. The results suggest that the program benefited a subset of elementary school students, but the effect on the subset was almost fully counterbalanced by no effects for the bottom 40 percentile. For middle-school students, the UQITT and ITT estimates are similar, pointing to a more homogeneous effect across the distribution in more advanced grades (see Figure \ref{fig:quantile} in the appendix). 

\begin{figure} [http]
		\caption{Unconditional Quantile ITT estimates for financial proficiency, elementary students - 95\%CI}
		\centering
		\includegraphics[scale=0.6]{Figure1.pdf}
	    \label{fig:te_firstcycle}
     
     \raggedright
        \fontsize{9}{10} \selectfont

     {Source: The financial proficiency is computed by CAEd based on student's answers on a standardized exam that used the Item Response Theory (IRT). The index is normalized so that the treatment effect can be measured in terms of standard deviations (SD). We first run this financial proficiency score on the treatment dummy, the six strata dummies, and on a constant. We then take the residual of this estimation and run a simultaneous quantile regression, for each one of the quantiles shown in the Figure, in which the dependent variable is the residual of the first regression and the independent variable is the treatment dummy. 1000 bootstrap replications.}
\end{figure}



%We find weaker effects concerning saving decisions though. The point estimate for the pooled sample is around half of the size of the consumption one and the analysis by grade level suggests that the effects on saving attitudes are statistically significant only for the fifth-graders. 

We further explore whether the program affected students' behavioral outcomes by looking at more objective measures of change in behavior such as the use of piggy banks, and investments in human capital. We expect that more patient students would be more likely to invest in human capital whose returns happen later in life. 

\autoref{tab:ITT_beha} shows impacts on students' likelihood to discuss financial subjects with their parents and friends, their likelihood to use a piggy bank (a proxy for actual savings), and their likelihood to receive a monthly allowance from parents (a proxy for students' access to informal sources of credit). Finally, we test whether the treatment increases children's access to financial services, such as debit and credit cards.\footnote{To measure students' access to financial services, we use a dummy variable that is one if a student has access to at least one service and zero otherwise.}  ITT estimates show that fifth graders are 5.7 pp (or 11.6 percent over the control mean) and 5.6 pp (or 28 percent over the control mean) more likely to talk to parents and friends about financial subjects respectively, and 2.7 pp. (or 71 percent over the control mean) more likely to use piggy banks (a proxy for actual savings). Overall, these results point to positive effects on awareness and saving outcomes. For middle school students, we find some indication that they become more likely to talk to friends about financial subjects. Unlike \cite{Bruhn2016}, we do not find any impact of the program on students' access to financial services. This is unsurprising in our context given that most students are young and do not deal with financial transactions in their daily lives.  



%\autoref{tab:implementation} shows that the percentage of students that had financial literacy classes during the whole school year is 10 percentage points higher among fifth-graders when compared to middle school students. As a consequence, the percentage of fifth-grade classes with more than 80\% of the textbooks covered is more than 10 percentage points higher than among middle school classes. 

Our estimates show that the program had larger impacts on financial proficiency in later grades and larger effects on attitudes and behavior in earlier grades. The results then suggest that the association between financial proficiency and behavioral change might not be straightforward. Individuals might have formal knowledge about different subjects, such as financial proficiency, but the application of the knowledge in real life can be discouraged by the institutional environment they are inserted (\cite{carpena2018causal}). To formally test the hypothesis that learning mediates change in behavior, we follow \cite{imai2011unpacking} and
decompose the treatment effect estimates into the average causal mediation effect (ACME), and
the average causal direct effect (ADE).\footnote {See the online appendix for a formal discussion of Casual Mediation Effects.} The goal of this exercise is to assess the share of the change in behavioral outcomes that could be attributed to improvements in financial proficiency (ACME) and the share that comes directly from the program (ADE).
The decomposition of ITT effects on ACME and ADE relies on the sequential ignorability assumption. The first part of the assumption assumes that conditioned on pre-treatment confounders the treatment assignment is orthogonal to the potential outcomes and the mediator outcome (financial proficiency). The second part of the assumption implies that the `observed mediator is ignorable given the actual treatment status and pre-treatment confounders \citep{imai2011unpacking}. This second part of the assumption is the strongest, even in the presence of random assignment. It poses that conditional on the actual treatment (not treatment assignment) and a vector of pre-treatment confounders, the outcome is independent of the mediator factor. Note that, in our case, this implies an assumption that attitudinal and behavioral outcomes are orthogonal to financial proficiency given the actual treatment and once we control for a vector of school characteristics. Since we did not design the experiment to obtain clean estimates of the ACME and ADE, we interpret our results as suggestive evidence.   

%To shed light on the link between financial proficiency and behavioral change, we estimate causal mediation effects, where we test whether an increase in the mediator variable, financial literacy, actually changes behavioral outcomes.  Since we did not design the experiment to test this hypothesis, 

%The model makes the strong assumption that the program impacts behavioral outcomes only indirectly.

%Finally, the proficiency test fails at capturing students' financial proficiency properly.

%While the ADE estimate refers to impacts on outcomes due to the treatment itself (everything else other than the mediator factor), the ACME estimate refers to the changes in outcomes that are mediated by the mediator factor, in our case, changes in financial proficiency. 




% Table generated by Excel2LaTeX from sheet 'Behaviroual outcomes'
\begin{table}[h!]
\begin{threeparttable}
\centering
\fontsize{9}{10} \selectfont
\caption{ITT estimates for behavioral outcomes}
\begin{tabular}{lccclccc}
\hline
\hline
           &     Pooled &    Elementary &     Middle &            &     Pooled &    Elementary &     Middle \\

           &            &   students &   students &            &            &   students &   students \\

       \multicolumn{ 4}{l}{{\bf Talk to parents}} &        \multicolumn{ 4}{l}{{\bf Talk to friends}} \\

 Treatment &      0.012 &   0.057 &     -0.012 &  Treatment &   0.035 &   0.056 &    0.023 \\

           &    (0.010) &    (0.019) &    (0.012) &            &    (0.009) &    (0.014) &    (0.011) \\

    pvalue &      0.216 &      0.003 &      0.306 &     pvalue &          0 &          0 &      0.044 \\

pvalue of RI &      0.222 &      0.004 &      0.345 & pvalue of RI &          0 &          0 &      0.058 \\

      Obs. &      10748 &       3875 &       6873 &       Obs. &      10725 &       3864 &       6861 \\

Number of clusters &        200 &        167 &        145 & Number of clusters &        200 &        167 &        145 \\

        R2 &      0.006 &      0.003 &      0.001 &         R2 &      0.008 &      0.012 &      0.005 \\

Strata F.E. &        Yes &        Yes &        Yes & Strata F.E. &        Yes &        Yes &        Yes \\

           &            &            &            &            &            &            &            \\
\hline
      \multicolumn{ 4}{l}{{\bf Piggy's bank use}} & \multicolumn{ 4}{l}{{\bf Use of finantial services}} \\

 Treatment &    0.022 &     0.027 &      0.017 &  Treatment &      0.002 &     -0.002 &      0.006 \\

           &    (0.011) &    (0.015) &    (0.013) &            &    (0.009) &    (0.012) &    (0.013) \\

    pvalue &      0.042 &      0.067 &      0.205 &     pvalue &      0.783 &       0.86 &      0.662 \\

pvalue of RI &      0.046 &      0.063 &      0.236 & pvalue of RI &      0.773 &      0.872 &      0.638 \\

      Obs. &      12739 &       5917 &       6822 &       Obs. &      12538 &       5718 &       6820 \\

Number of clusters &        201 &        172 &        145 & Number of clusters &        201 &        172 &        145 \\

        R2 &      0.028 &      0.038 &       0.01 &         R2 &      0.005 &      0.004 &      0.006 \\

Strata F.E. &        Yes &        Yes &        Yes & Strata F.E. &        Yes &        Yes &        Yes \\

           &            &            &            &            &            &            &            \\
\hline
             \multicolumn{ 4}{l}{{\bf Allowance}} &     {\bf } &            &            &            \\

 Treatment &      0.013 &      0.017 &       0.01 &            &            &            &            \\

           &    (0.008) &    (0.014) &    (0.010) &            &            &            &            \\

    pvalue &       0.12 &      0.222 &      0.303 &            &            &            &            \\

pvalue of RI &      0.111 &      0.245 &      0.305 &            &            &            &            \\

      Obs. &      10699 &       3862 &       6837 &            &            &            &            \\

Number of clusters &        200 &        167 &        145 &            &            &            &            \\

        R2 &      0.009 &      0.012 &      0.007 &            &            &            &            \\

Strata F.E. &        Yes &        Yes &        Yes &            &            &            &            \\
\hline
\end{tabular}  
\label{tab:ITT_beha}
\begin{tablenotes}
    \fontsize{9}{9}\selectfont 
    \item Note: Clustered standard errors at the school level. RI stands for Randomization Inference. \\
    The dependent variables are equal to 1 if students talk to parents or friends about financial subjects, if they use a piggy bank, financial services, or receive an allowance from their parents, and 0 otherwise. 
\end{tablenotes} 
\end{threeparttable}
\end{table}%



%The association between changes in financial proficiency and behavioral outcomes is much weaker. 

According to the estimates in \autoref{tab:acme}, while we find some evidence that financial proficiency mediates the improvement of middle students' hypothetical attitudes in scenarios involving savings and consumption, it does not mediate changes in their behavior. For elementary school students, all of the effects seem to be direct, i.e., not mediated by changes in financial proficiency (\autoref{tab:acme}. In other words, the ADE and ACME estimates suggest that financial education programs affect students' attitudes and behavior through pathways other than the accumulation of formal knowledge. Overall, the estimates do not support the hypothesis that changes in financial proficiency are a necessary condition to change individuals' behavioral outcomes, a finding consistent with \cite{carpena2018causal}, and \cite{de2021effect}.  

%changes in financial proficiency seems to mediate changes in children's attitudes towards hypothetical consumption and saving decisions. Notwithstanding, this association seems to hold only for middle school grades. 

To assess whether the results on both financial proficiency, attitudes, and behavioral outcomes are picking up some heterogeneity in the data, we carried out a series of subgroup analyses. First, we compared the effects on boys and girls. The point estimates are larger for boys, but the difference in estimates between boys and girls is not statistically significant in most cases. Second, we tested whether the impacts varied by poverty status. We aimed to assess whether poorer kids - those participating in the Brazilian conditional cash transfer program - would respond less to an intervention intended to help kids become more forward-looking. In fact, the estimates on financial proficiency and consumption attitudes are larger in absolute terms for the poor, but the differences in coefficients between poor and non-poor are not statistically significant for the most part. The heterogeneous effect analyses are available in the online appendix.   

\subsection{Robustness Check} 

Back in 2015, we did not have access to administrative student-level data to check the balance between treatment and control groups and investigate the impacts of the program on students' performance in reading and math. As described in \autoref{sec:data_c}, in 2020 and 2021 INEP granted us access to its facilities to complement our analyses using student-level data. These data allow us to rerun the balance tables at the student level and, more importantly, look at the program's impacts on both contemporary and longer-term human capital outcomes.

As discussed in \autoref{sec:pilot}, school principals' choice of the classrooms to participate in the pilot might not have been random. Consequently, if school principals strategically picked the best classrooms aiming to show the benefits of the program, or the worst classrooms betting on a potential complementarity between financial education and other school subjects, our estimates would be biased as we could be comparing treated and control students with a different pre-treatment level of skills. To assess students' pre-treatment level of proficiency in Portuguese and math, we look at their scores in the Brazilian standardized exam when in grade 5. Indeed, the average test scores in grade 5 for students attending the seventh and ninth grades in 2015 were higher in the control group (\autoref{tab:balance_adm2}). To check whether this imbalance in pre-treatment skill level affects our results, we re-estimate \autoref{eq:Y} for seventh and ninth graders controlling for their pre-treatment math and reading test scores.\footnote{Seventh and ninth-graders are the ones for which the pre-treatment performance in reading and math is available. For seventh-graders, we have their fifth-grade performance in reading and math in 2013. For ninth-graders, we have their fifth-grade performance in 2011.} 

%As described in \autoref{sec:data_c}, we gathered pre-treatment student performance in reading and math by merging pilot students with administrative INEP datasets. Among seventh and ninth-graders who finished the financial proficiency assessment, we found around two-thirds in the 2013 and 2011 waves of \textit{Prova Brasil}, when they were supposedly enrolled in fifth grade (assuming that they did not repeat any grade until 2015). The percentage is balanced between the treatment and control groups.

Controlling for past performance in math and reading leads to a substantial gain in precision. Among seventh graders, we observe a slight increase in the point estimate for financial proficiency scores (from 0.09 SD to 0.098 SD). For ninth-graders, we find a meaningful increase in financial proficiency scores, with the point estimate jumping from 0.11 SD to 0.15 SD (\autoref{tab:add_5prof_control}), though the difference is not statistically significant. The updated estimates indicate that we would underestimate the impacts of the pilot on financial proficiency if we did not control for an imbalance in pre-treatment proficiency in math and reading at the student level, particularly for ninth graders.  Controlling for past performance in math and reading scores does not change the point estimates for consumption and saving attitudes. 


\subsection{Implementation and IV Results}

The following analysis seeks to shed some light upon the role played by the quality of the program's implementation on its impacts. This exercise is useful in that it avoids hasty conclusions that a given intervention did not have the expected effect when in fact the program was not carried out as planned. 
We use a variable that informs the percentage of the syllabus (use of the teaching material) covered during the school year. As discussed earlier, the material was delivered somewhat late at the treated schools, and that eventually hindered the exposure to the program and consequently impacted the intensity of treatment. Approximately one-third of the teachers claimed they covered up to 60\% of the syllabus. 
We first carry out a correlation analysis testing whether the effects are associated with the quality of the program's implementation. We do so by splitting the treatment group into two groups: schools with up to 60\% of the syllabus covered, and schools with more than 60\% of the syllabus covered.

Since the quality of the program's implementation is endogenous, we follow \cite{frisancho2020impact} and use the random assignment to treatment as an instrumental variable for the program's implementation.
The IV estimates provide the local average treatment effect (LATE) of the program by comparing the subgroup of treated schools that received the treatment (supervisor/teacher training and teachers who used the textbooks) with the control schools. We create a binary variable, $Z_$, that is 1 for schools with at least 60\% of the syllabus covered during the academic year, and 0 otherwise leveraging information gathered from teachers and students on whether the textbooks were incorporated into the classroom during the school year.\footnote{Unfortunately, the variable is categorical and we are unable to define the dummy at the median of the distribution or to estimate marginal treatment effects.} Since the data show that some students and teachers from control schools apparently received and used the textbooks, we use a two-stage least squares to estimate the LATE:


First-Stage: \vspace{-50pt}

\begin{equation}
     Z_{iksm}= \gamma_{1} + \gamma_{2} T_{hksm} + \Sigma{C_j} + \omega_{1iksm},  h = 1, 2. 
     \label{eq:Y2}
\end{equation}

Second-Stage:\vspace{-50pt}

\begin{equation}
    Y_{iksm}= \lambda_{1} + \lambda_{2} \widehat{Z_{iksm}} + \Sigma{C_j} + \omega_{2iksm} 
    \label{eq:Y3} 
\end{equation}

The parameter of interest, $\lambda_2$, is the LATE. The IV estimates are shown in Table \ref{tab:iv}. The estimates 
are slightly larger and less precise than ITT ones, suggesting that the program did not have large effects even when better implemented.   
\vspace{10 pt}

\begin{table}[ht]
\centering
\begin{threeparttable}
\fontsize{10}{11}\selectfont
\caption{IV estimates}
\begin{tabular}{lcccccc}
\hline
\hline
           &   \multicolumn{ 2}{c}{} &   \multicolumn{ 2}{c}{} &   \multicolumn{ 2}{c}{} \\

           & \multicolumn{ 2}{c}{Pooled} & \multicolumn{ 2}{c}{Elementary} & \multicolumn{ 2}{c}{Middle school} \\
\hline
{\bf Financial Proficiency} &        ITT &       LATE &        ITT &       LATE &        ITT &       LATE \\
           &   \multicolumn{ 2}{c}{} &   \multicolumn{ 2}{c}{} &   \multicolumn{ 2}{c}{} \\

  Estimate &            &            &            &            &            &            \\

           &    0.073 &    0.085 &      0.047 &      0.055 &    0.099 &    0.115 \\

    N. obs &    (0.034) &    (0.039) &    (0.046) &    (0.054) &    (0.042) &    (0.049) \\

        R2 &     14,655 &     14,655 &      7,641 &      7,641 &      7,014 &      7,014 \\

           &      0.167 &      0.166 &       0.15 &      0.149 &      0.188 &      0.189 \\
\hline
{\bf Consumption Index} &            &            &            &            &            &            \\

  Estimate &   0.094 &   0.110 &   0.122 &   0.146 &   0.082 &   0.095 \\

           &    (0.026) &    (0.031) &    (0.042) &    (0.050) &    (0.031) &    (0.036) \\

    N. obs &     10,220 &     10,220 &      3,648 &      3,648 &      6,572 &      6,572 \\

        R2 &      0.089 &      0.087 &      0.119 &      0.117 &      0.076 &      0.076 \\

           &            &            &            &            &            &            \\
\hline
{\bf Savings Index} &            &            &            &            &            &            \\

  Estimate &    0.055 &    0.064 &    0.087 &    0.104 &      0.044 &      0.051 \\

           &    (0.025) &    (0.029) &    (0.041) &    (0.049) &    (0.029) &    (0.034) \\

    N. obs &     10,022 &     10,022 &      3,549 &      3,549 &      6,473 &      6,473 \\

        R2 &      0.046 &      0.046 &      0.096 &      0.095 &      0.025 &      0.025 \\

\hline
\label{tab:iv}
\end{tabular}  
 \begin{tablenotes}
    \fontsize{9}{9}\selectfont 
    \item Note: Clustered standard errors at the school level. Students are considered treated if their teachers answered that at least part of the content of financial education was covered in class. For 8.5\% of the treated students and 7.5\% of the control students, we do not have a teacher assigned. In these cases, we cannot infer actual treatment using teachers' answers on the usage of the material. We then consider a student as treated if at least 50\% of her peers in the classroom state they received the financial literacy textbook.  \\      
\end{tablenotes}
\end{threeparttable}
\end{table}


\subsection{Longer-term Effects on Behavioral Outcomes: Investment in Human Capital} \label{subsec:human_capital}

In this section, we assess whether the pilot impacts students' human capital investment decisions. The hypothesis is that if students understand the importance of education for their future earnings and become more patient and forward-looking following the pilot, they would be more likely to study harder, potentially improve their academic performance in reading and math, and be less likely to drop out of school (\cite{Bruhn2016}).  This has a direct policy implication as it would indicate whether offering financial education to elementary and middle school students could improve individuals' outcomes other than their financial proficiency with potential benefits to their financial health and labor market outcomes in the longer term.

%inter-temporal decision-making
%Also, an increase in financial knowledge could eventually lead to higher math and reading skills if these are complementary skills as suggested by figure \ref{fig:prof_ideb}.

We start by checking if the program has an immediate impact on the performance of fifth and ninth-graders in math and reading in 2015, the year the financial education program was piloted.\footnote{Fifth and ninth grades are the ones for which the performance in math and reading is available in 2015 due to their participation in \textit{Prova Brasil}. See \autoref{tab:inep_data}.} Figure \autoref{fig:prof_ideb} shows a significant correlation between students' financial proficiency score and their reading and math performance. We then investigate whether there is any causal relationship in place. Our results do not point to a contemporary effect of financial education on either math or reading skills for both grades (\autoref{tab:long-lasting}). Although the nil effects on math and reading proficiency of fifth graders could be explained by the fact that the intervention had a small effect on a subset of the students as suggested by the UQITT estimates, the lack of treatment effects for grade 9 is less expected. 

%Also, the lack of impact could be due to implementation issues. Even though the program's implementation was slightly be better in elementary schools grades, 60\% of fifth-graders had access to the financial literacy textbooks either only in the first half or only in the second half of the school year (\autoref{tab:implementation}). As consequence, only a small percentage (one-third) of fifth-graders had more than 80\% of the curricula covered in the academic year. Students exposed to the lessons only in the second semester (almost 40\%) had just three months between the introduction of the lessons and the application of \textit{Prova Brasil}.\footnote{Taking into  consideration that the beginning of the second semester is in August, and the application of \textit{Prova Brasil} is at the end of October or at the beginning of November.} 

These results seem to suggest that teachers involved in this pilot evaluation did not prioritize financial education at the expense of other school subjects (e.g., teach to the test), a concern the research team and some school principals had back in the day. Overall, our estimates are in line with \cite{batty2020experiential} that also finds evidence of no substitution from the core school curriculum. 

%Notice that these results somehow weaken concerns with teachers' strategic behavior such as 

%However, as with fifth graders, students in grade 9 faced various challenges with program's implementation. For instance, one fifth of the students had more than 80\% of the syllabus covered. More than 50\% of ninth graders had their financial education lessons concentrated in the second half of the school year, which led to only three months of financial proficiency exposure before the proficiency test in math and reading (\textit{Prova Brasil}).  




Besides the contemporary effects on learning outcomes, the data allow us to test whether the program impacted students' human capital accumulation. We first assess whether the pilot affected students' probability of being retained or dropping out of school between 2015 and 2018. Second, for third and seventh graders, we assess their performance in math and reading two years after the intervention (2017) when they reached fifth and ninth grades, respectively.\footnote{For those who did not repeat between third and fifth grades and between seventh and ninth graders, respectively.} Finally, for ninth-graders, we assess whether the intervention affects their probability of finishing high school, taking the ENEM test, and the test takers' ENEM proficiency score. 

We do not find any indication that the financial education program affects students' academic achievements -- test scores and school progression -- in the longer term (\autoref{tab:long-lasting}). We interpret these findings as evidence that the program does not make students become more forward-looking with regard to their human capital accumulation decisions, which might be associated with poor program implementation as described in \autoref{sec:implementation}.

\begin{landscape}
\newgeometry{left = 1.8cm, right = -10cm, top = 0.5cm, bottom = -2cm}

\begin{table}[ht]
\centering
\begin{threeparttable}
\fontsize{9}{10}\selectfont
\caption{ITT estimates for Human Capital Accumulation Outcomes}
% Table generated by Excel2LaTeX from sheet 'Table18'
\begin{tabular}{lccccccccccccc}
\hline
\hline
           & \multicolumn{ 2}{c}{3rd grade} &            & \multicolumn{ 2}{c}{5th grade} &                    \multicolumn{ 4}{c}{7th grade} &            &       \multicolumn{ 3}{c}{9th grade} \\

           &       Only & Strata FE  &            &       Only & Strata FE  &            &       Only & Strata FE  & Strata FE +  &            &       Only & Strata FE  & Strata FE +  \\

           &     strata &  Financial &            &     strata &  Financial &            &     strata &  Financial & Performance &            &     strata &  Financial & Performance \\

           &            &  proficiency &            &            &  proficiency &            &            &  proficiency &  5th grade &            &            &  proficiency &  5th grade \\
\hline
{\bf Repetition} &      0.017 &      0.010 &            &    0.029 &     0.025 &            &     -0.007 &     -0.002 &     -0.008 &            &      0.016 &     0.031 &      0.018 \\

           &    (0.017) &    (0.019) &            &    (0.014) &    (0.015) &            &    (0.014) &    (0.015) &    (0.016) &            &    (0.016) &    (0.017) &    (0.018) \\

           &            &            &            &            &            &            &            &            &            &            &            &            &            \\

{\bf Dropout} &      0.005 &      0.007 &            &      0.011 &      0.000 &            &      0.006 &      0.002 &      0.002 &            &      0.033 &      0.028 &      0.020 \\

           &    (0.006) &    (0.006) &            &    (0.010) &    (0.008) &            &    (0.018) &    (0.014) &    (0.012) &            &    (0.020) &    (0.018) &    (0.019) \\

           &            &            &            &            &            &            &            &            &            &            &            &            &            \\

{\bf Portuguese} &     -0.029 &     -0.053 &            &     -0.017 &     -0.032 &            &     -0.004 &     -0.057 &     -0.005 &            &     -0.015 &    -0.063 &      0.028 \\

{\bf Prova Brasil} &    (0.047) &    (0.051) &            &    (0.046) &    (0.031) &            &    (0.047) &    (0.041) &    (0.034) &            &    (0.042) &    (0.036) &    (0.036) \\

           &            &            &            &            &            &            &            &            &            &            &            &            &            \\

{\bf Math} &     -0.066 &    -0.096 &            &     -0.033 &     -0.045 &            &     -0.005 &     -0.060 &     -0.009 &            &    -0.077 &   -0.094 &     -0.020 \\

{\bf Prova Brasil} &     (0.053) &     (0.055) &            &    (0.049) &    (0.039) &            &    (0.051) &    (0.048) &    (0.038) &            &    (0.042) &    (0.037) &    (0.039) \\

           &            &            &            &            &            &            &            &            &            &            &            &            &            \\

{\bf Finish} &            &            &            &            &            &            &            &            &            &            &     -0.035 &     -0.035 &     -0.019 \\

{\bf high school} &            &            &            &            &            &            &            &            &            &            &    (0.025) &    (0.024) &    (0.023) \\

           &            &            &            &            &            &            &            &            &            &            &            &            &            \\

{\bf Participation} &            &            &            &            &            &            &            &            &            &            &     -0.025 &     -0.016 &     -0.012 \\

{\bf in ENEM} &            &            &            &            &            &            &            &            &            &            &    (0.024) &    (0.022) &    (0.023) \\

           &            &            &            &            &            &            &            &            &            &            &            &            &            \\

{\bf Portuguese} &            &            &            &            &            &            &            &            &            &            &     -0.046 &     -0.004 &     -0.025 \\

{\bf ENEM} &            &            &            &            &            &            &            &            &            &            &     (0.036) &      (0.05) &     (0.031) \\

    {\bf } &            &            &            &            &            &            &            &            &            &            &            &            &            \\

{\bf Math} &            &            &            &            &            &            &            &            &            &            &     -0.028 &     -0.035 &      0.002 \\

{\bf ENEM} &            &            &            &            &            &            &            &            &            &            &     (0.047) &     (0.063) &     (0.045) \\

      Obs. &       4650 &       3635 &            &       4755 &       4006 &       4697 &            &       3640 &       3137 &            &       4455 &       3374 &       2851 \\
\hline
\hline 
\label{tab:long-lasting}
\end{tabular}  
 \begin{tablenotes}
    \fontsize{9}{9}\selectfont 
    \item Note: Clustered standard errors at the school level. \\
    Grade retention assumes a value of 1 if the student is retained in the same grade at the end of the school year, and 0 otherwise. Dropout is equal to 1 if the student drops out at the end of the school year, and 0 otherwise. Finish high school is equal to 1 if the student finishes high school, and 0 otherwise. Participating in ENEM is equal to 1 if the students participate in the ENEM assessment, and 0 otherwise. Portuguese and Math are the proficiency scores in \textit{Prova Brasil} on the SAEB scale. 
      
\end{tablenotes}
\end{threeparttable}
\end{table}
\end{landscape}